% figures/ch09/fig02-final-layout.tex — auto-extracted from chapter source
\begin{tikzpicture}[
  block/.style={draw, minimum width=7cm, font=\small, anchor=south west},
  addr/.style={font=\ttfamily\scriptsize\color{gray}, anchor=east},
  arr/.style={-{Stealth[length=3pt]}, thick, blue!60!black},
]
  % 低地址：未映射
  \node[block, minimum height=2.0cm, fill=gray!10, dashed] (unmapped) at (0, 0)
    {\parbox{6.5cm}{\centering\small\textit{未映射（将来的用户空间）}\\
     访问触发 \#PF}};

  % 直接映射区
  \node[block, minimum height=1.5cm, fill=green!10] (direct) at (0, 2.0)
    {\parbox{6.5cm}{\centering\small\textbf{直接映射区}\\
     $\text{virt} = \text{phys} + \hex{C0000000}$}};

  % 动态分配区
  \node[block, minimum height=0.8cm, fill=yellow!10] (alloc) at (0, 3.5)
    {\parbox{6.5cm}{\centering\small 动态分配区（vmm\_alloc\_pages）}};

  % 内核堆
  \node[block, minimum height=0.8cm, fill=orange!10] (heap) at (0, 4.3)
    {\parbox{6.5cm}{\centering\small 内核堆（\hex{E0000000}+）}};

  % 地址标注
  \node[addr] at (-0.2, 0)    {\hex{00000000}};
  \node[addr] at (-0.2, 2.0)  {\hex{C0000000}};
  \node[addr] at (-0.2, 3.5)  {direct\_map\_end};
  \node[addr] at (-0.2, 4.3)  {\hex{E0000000}};
  \node[addr] at (-0.2, 5.1)  {\hex{FFFFFFFF}};

  % 物理内存箭头
  \node[draw, minimum width=2.5cm, minimum height=1.5cm, fill=blue!8,
        font=\small, anchor=west] (phys) at (8.5, 2.0)
    {\parbox{2cm}{\centering 物理\\内存}};
  \draw[arr] (7.0, 2.75) -- (phys.west) node[midway, above, font=\scriptsize] {1:1};
\end{tikzpicture}
